\subsubsection{Pipelines}

% Original text from ms.tex
% \spi\ pipeline optimizations, updates, and testing have been managed by part of this team since 2018 (principally at UdeM, IRAP, IPAG, and LNA). Planned improvements for \wena\ include the implementation of line-by-line RV extraction and enhanced telluric calibration for \esp\ (building on \citealt{artigau_line-by-line_2022}), and homogeneous spectropolarimetric processing and distribution via PolarBase, which already archives \esp\ and \spi\ data.

The team has direct, long-term experience in the development, validation, and operation of the reduction and analysis pipelines required for \wena. On the \spi\ side, \texttt{APERO} \citep{cook_apero_2022} has been built and maintained by members of the collaboration since first light in 2018, with contributions led principally from UdeM, IRAP, IPAG, and LNA. \texttt{APERO} handles the full reduction chain from raw detector frames to extracted spectra, wavelength calibration, blaze correction, telluric correction, polarimetric products, and precision-RV measurements, and it will remain the backbone of the \spi\ data flow throughout \amakihi.

For \wena, we will extend this expertise to a homogeneous processing framework across \spi\ and \esp. Priority developments include the adaptation and continued improvement of the \esp\ reduction for the \wena\ observing mode, the implementation of line-by-line RV extraction on \esp\ data and on joint time-series products \citep{artigau_line-by-line_2022}, and improved telluric handling over the full spectral range. These core reductions will be complemented by advanced analysis layers already developed within the team, including $dTemp$ for tracing subtle stellar-temperature variations \citep{artigau_measuring_2024}, WAPITI for posterior mitigation of residual telluric and instrumental systematics in RV time series \citep{ould-elhkim_wapiti_2023}, and chromatic activity diagnostics tailored to disentangle stellar variability from planetary signals \citep{larue2025}.

The same effort will ensure homogeneous spectropolarimetric products across the survey. Least-Squares Deconvolution and related magnetic diagnostics will be produced within the established \esp/\spi\ analysis ecosystem, and the resulting reduced and value-added products will be distributed through channels already familiar to the community: science-ready reduced data through CADC and spectropolarimetric products through PolarBase, which already archives \esp\ and \spi\ observations. In this way, the pipeline strategy for \amakihi\ is not a speculative development but a continuation and consolidation of tools that have already delivered mature science results on both the RV and spectropolarimetric fronts.